
Training data deduplication is commonly applied to improve language model pretraining quality. When Pythia models trained on the Pile are compared to models trained on its deduplicated variant (dedup-Pile) at token-count-matched checkpoints, deduplication does not produce a uniform accuracy change: MMLU accuracy decreases significantly while HellaSwag and ARC-Challenge accuracy increases. This divergence is consistent with a contamination-correction interpretation: the per-benchmark accuracy differential (dedup-Pile minus Pile) correlates positively with each benchmark's estimated 13-gram n-gram overlap with the Pile training corpus (Pearson r = 0.632, p = 0.0086, n = 16 observations across 4 benchmarks and 4 model sizes from 160M to 6.9B parameters). The benchmark with the highest literature-estimated contamination in Pile (MMLU, 5.5\% 13-gram overlap) shows a Bonferroni-corrected significant accuracy reduction under deduplication (t = -5.574, p = 0.0114), while the lowest-contamination benchmark (WinoGrande, 2.5\%) is negligibly affected. Additionally, token-count matching --- identifying Pile checkpoints at equivalent total token counts rather than equivalent training steps --- recovers a Delta-r = +0.093 stronger contamination signal than step-matching. The near-memorization mechanism hypothesized to explain this signature (as measured by min-k\% probability scores at Pythia-1B scale) was not confirmed: dedup-Pile models showed higher min-k\% scores than Pile models on all four benchmarks, opposite to prediction. This negative finding is reported as a limitation. These results reframe deduplication's benchmark-level effects as contamination correction rather than uniform quality improvement.

